\documentclass[10pt,a4paper]{amsart}
\usepackage[margin=19mm]{geometry}
\usepackage{amsmath,amssymb,mathtools}
\usepackage[scr=rsfs,cal=euler]{mathalfa}
\usepackage{xfrac}
\usepackage[unicode,hidelinks,bookmarks=false]{hyperref}
\newcommand{\ie}{i.e.\ }
\newcommand{\cf}{cf.\ }
\newcommand{\resp}{respectively\ }
\newcommand{\Subsection}{\subsection}
\providecommand{\nobreakdash}{\nobreak\hbox{-}}
\setlength{\emergencystretch}{2em}
\multlinegap0pt
\usepackage{amssymb}
\newtheorem{theorem}{Theorem}
\newcommand{\Tr}[0]{\operatorname{Tr}}
\title{Tracing the Loop: \\ \large Non-Causal Computation, Partial Traces, \& Postselected Entanglement}
\author{Mark Carney}
\date{Source: arXiv:2609.08716v2 (CC BY 4.0)}
\begin{document}
\maketitle

\noindent\emph{Source and licence.} Tracing the Loop: \\ \large Non-Causal Computation, Partial Traces, \& Postselected Entanglement. Version arXiv:2609.08716v2 (CC BY 4.0), distributed by arXiv under the Creative Commons Attribution 4.0 International licence, http://creativecommons.org/licenses/by/4.0/, as recorded in the arXiv licence metadata for this exact version and shown on the abstract page https://arxiv.org/abs/2609.08716v2. Full licence notice: \texttt{licenses/arxiv-2609.08716-CC-BY-4.0.txt}.\par\smallskip

\noindent\textit{Editorial scope.} Counted unit: the article's theorem theorem (source0.tex lines 171--177). The article's complete original proof of it is delivered with it (source0.tex lines 179--195, one \texttt{proof} environment of 17 lines). No mathematics is written, repaired or re-derived: the statement and the proof are the article's own lines, in the article's own notation, and no retained line is substituted or re-flowed.  Retained uncounted as the source-local context the proof needs: lines 108--110.  Nothing was omitted from any retained span, so the package carries no omission record.  No retained line contains a citation, so the wrapper carries no bibliography and no bibliography entry.  No retained line contains a figure environment, an image-inclusion command or drawing code, so no image is delivered and the package declares no asset dependency.  Editorial formatting only, in addition: an AMS \texttt{amsart} preamble that restates the article's own declarations and macros used by the retained lines, namely the environments \texttt{theorem} on separate counters, together with the standard packages those lines need.  The retained environments print in this document as Theorem 1 in order of appearance, whereas the article prints the same items with the numbering of its own full document.  The complete native source of the article as distributed in the arXiv e-print for this exact version is bundled unchanged as \texttt{sources/209739/source0.tex}.
\begin{align}\label{eq:morphismP}
    P: A \times C \rightarrow X \times C
\end{align}

\begin{theorem}
    Let $P$ be deterministic with the same definition as eq. (\ref{eq:morphismP}), so that there exist component functions $h: A \times C \to X$ and $g: A \times C \to C$ such that $$ P(x, c_{\mathrm{out}} \mid a, c_{\mathrm{in}}) = \delta_{x, h(a, c_{\mathrm{in}})} \delta_{c_{\mathrm{out}}, g(a, c_{\mathrm{in}})} $$ The following are equivalent:
    \begin{enumerate}
        \item $\Tr_C(P) \in \mathbf{Stoch}$
        \item For each $a$ there exists a unique $c_a$ such that $g(a, c_a) = c_a$.
    \end{enumerate}
\end{theorem}

\begin{proof}
    Substituting the categorical trace formula we get 
    \begin{align*}
        \Tr_C(P)(x \mid a) &= \sum_{c \in C} P(x, c \mid a, c) \\
        &= \sum_{c \in C} \delta_{x, h(a, c)} \delta_{c, g(a, c)}
    \end{align*}

    The factor $\delta_{c, g(a, c)} = 1$ precisely when $g(a,c) = c$, therefore $$ \Tr_C(P)(x \mid a) = |F| $$ where $$ F = \{ c \in C : g(a,c)=c \text{ and } h(a,c) = x \}$$ Summing over $x$:
    \begin{align*}
        \sum_{x \in X} \Tr_C(P)(x\mid a) &= \sum_{x \in X} \sum_{c \in C} \delta_{x, h(a, c)} \delta_{c, g(a, c)} \\
        &= \sum_{c \in C} \delta_{c, g(a,c)} \sum_{x\in X} \delta_{x, h(a, c)} \\
        &= \sum_{c \in C} \delta_{c, g(a,c)}
    \end{align*}
    The last step following from the fact that $\sum_{x\in X} \delta_{x, h(a, c)} = 1$ for every fixed $(a,c)$. The right hand side is precisely the fixed points of $g$, and so $\Tr_C(P) \in \mathbf{Stoch}$ if and only if $g$ has exactly one fixed point for each $a$. This completes the equivalence.
    
    Let $c_a$ be the fixed point for some $a$ such that $g(a,c_a) = c_a$, then every term in the sum $$ \sum_{x \in X} \Tr_C(P)(x\mid a) = \sum_{x \in X} \sum_{c \in C} \delta_{x, h(a, c)} \delta_{c, g(a, c)} $$ goes to zero except when $c = c_a$. So the sum reduces to the single contribution from $$ \Tr_C(P)(x \mid a) = \delta_{x, h(a, c_a)} $$ Thus the traced process is the deterministic function $$ a \mapsto h(a, c_a)$$
\end{proof}

\end{document}
